%% ----------------------------------------------------------------------------
%%                              Chinese Abstract
%% ----------------------------------------------------------------------------
\begin{abstract}{图推理, 大语言模型, 多智能体}
图推理面向图结构数据中的关系、路径、连通性和组合约束进行计算，是图数据分析和智能决策中的基础任务，在路径规划、网络分析、组合优化以及知识推断等场景中均有应用。随着大语言模型在语义理解、步骤化推理和代码生成方面能力的提升，利用大语言模型自动求解图推理问题成为一种可行思路。但与一般文本推理不同，图推理答案必须严格满足节点、边、路径、连通关系及任务约束。现有端到端方法通常先将图结构转写为文本，再由模型直接生成推理过程或答案。在图规模较大、约束条件较多或输入序列较长时，这种方式容易遗漏局部拓扑关系，进而产生结构理解偏差或看似合理但不满足约束的答案。另一类方法借助历史任务样本、经验代码库或专门微调模型来生成程序，能够在部分固定任务上改善结果，但其效果依赖已有样例和适配数据，数据构建、维护与更新成本较高；当任务形式变化、外部知识更新或跨基准迁移时，泛化能力仍受到限制。

针对上述问题，本文提出基于规划反馈的多智能体图推理协同求解框架 PlanGraph。该框架将图推理过程拆分为规划生成、规划引导检索、执行验证与反馈修复等相互衔接的阶段。具体而言，问题智能体和规划智能体首先对自然语言问题、图结构描述和输出要求进行解析，并生成伪代码形式的规划结果，用于显式记录任务目标、算法步骤和关键约束。随后，检索智能体不再直接依据原始问题检索，而是利用规划结果从图算法文档知识库中获取相关 API、算法说明和实现知识，并在检索结果不足时进行片段精炼或补充检索。编码智能体根据规划与检索结果生成可执行程序，再通过运行测试、错误反馈、推理修正和有界重生成形成验证闭环。通过这种设计，规划在前向求解过程中约束知识调用和代码生成，反馈则在反向过程中定位并修正实现偏差、约束遗漏和输出不一致问题，从而提高程序化图推理的稳定性、正确性和可解释性。

本文在 GraphArena、NLGraph、GraphWiz、LLM4DyG 和 GraphInstruct 五个公开图推理基准测试集上对所提方法进行了系统评估。PlanGraph 在五个基准测试集中有四个取得最优结果，平均精确匹配率达到 97.66\%，总体结果优于现有方法，并在经典图算法、组合优化和指令式图推理任务中表现出较强的鲁棒性。实验结果显示，将显式规划、外部知识检索与程序验证组织为反馈闭环，能够有效缓解大语言模型在图推理任务中的幻觉与不稳定问题，并为大语言模型与图算法工具协同求解复杂结构化问题提供一条具有工程可行性和可解释性的实现路径。
\end{abstract}

%% ----------------------------------------------------------------------------
%%                              English Abstract
%% ----------------------------------------------------------------------------
\begin{englishabstract}{graph reasoning, large language models, multi-agent systems}
Graph reasoning computes relations, paths, connectivity, and combinatorial constraints in graph-structured data. It is a fundamental task in graph data analysis and intelligent decision-making, and has been applied to path planning, network analysis, combinatorial optimization, and knowledge inference. With the improvement of large language models in semantic understanding, step-by-step reasoning, and code generation, using large language models to automatically solve graph reasoning problems has become a feasible direction. However, unlike general textual reasoning, answers to graph reasoning tasks must strictly satisfy node, edge, path, connectivity, and task-specific constraints. Existing end-to-end methods usually first rewrite graph structures as text, and then ask the model to directly generate the reasoning process or the final answer. When the graph scale is large, the constraints are complex, or the input sequence is long, this approach may miss local topological relations, leading to biased structural understanding or seemingly plausible answers that do not satisfy the constraints. Another line of methods generates programs with the help of historical task samples, empirical code repositories, or specially fine-tuned models. These methods can improve results on some fixed tasks, but their effectiveness depends on existing examples and adaptation data, which brings relatively high costs for data construction, maintenance, and updating. When task forms change, external knowledge is updated, or benchmarks are transferred, their generalization ability remains limited.

To address these problems, this thesis proposes PlanGraph, a planning-feedback-based multi-agent collaborative framework for graph reasoning. The framework decomposes the graph reasoning process into connected stages, including planning generation, planning-guided retrieval, execution verification, and feedback repair. Specifically, the problem agent and the planning agent first parse the natural language question, graph structure description, and output requirements, and generate a pseudocode-style plan to explicitly record the task objective, algorithmic steps, and key constraints. Then, instead of retrieving directly based on the original question, the retrieval agent uses the plan to obtain relevant APIs, algorithm descriptions, and implementation knowledge from a graph algorithm documentation knowledge base, and refines knowledge fragments or performs supplementary retrieval when the retrieved results are insufficient. The coding agent generates executable programs based on the plan and retrieved results, and further forms a verification loop through execution tests, error feedback, reasoning-based correction, and bounded regeneration. Through this design, planning constrains knowledge invocation and code generation in the forward solving process, while feedback locates and corrects implementation deviations, missing constraints, and output inconsistencies in the reverse process, thereby improving the stability, correctness, and interpretability of programmatic graph reasoning.

This thesis systematically evaluates the proposed method on five public graph reasoning benchmarks: GraphArena, NLGraph, GraphWiz, LLM4DyG, and GraphInstruct. PlanGraph achieves the best performance on four of the five benchmarks, with an average exact match rate of 97.66\%. Overall, it outperforms existing methods and shows strong robustness in classical graph algorithms, combinatorial optimization, and instruction-style graph reasoning tasks. These results demonstrate that organizing explicit planning, external knowledge retrieval, and program verification into a feedback loop can effectively alleviate hallucination and instability in LLM-based graph reasoning tasks, and provide an engineering-feasible and interpretable implementation path for collaborative solving of complex structured problems with LLMs and graph algorithm tools.
\end{englishabstract}
